%% ARXIV-ONLY fork of sections/calculus.tex — arxiv.tex inputs this file;
%% main.tex keeps sections/calculus.tex untouched (the POPL submission is
%% frozen). In this version the DIRECTIONAL square is the primary notion
%% and the exact square is its identity-embedding special case; the POPL
%% version presents the exact square only, and its former \ifarxiv
%% directional subsection is absorbed here. Cherry-pick toward a future
%% POPL revision by diffing against sections/calculus.tex.
%%
%% NUMBERING CONTRACT: items 3.1–3.9 must keep their numbers — the
%% appendix (body-arxiv.tex, and the frozen standalone body.tex) and the
%% published arXiv v1 refer to them. New numbered material goes after
%% 3.9 only. Current: 3.10 = prop:exactid, 3.11 = prop:lax (its v1
%% number, preserved).

\section{A calculus of pairs}
\label{sec:calculus}

This section makes the informal contract of \Cref{sec:overview} precise:
languages, behaviors, and projections (\S\ref{sec:languages}); pairs and
their directional commuting squares (\S\ref{sec:pairs}); the pasting
theorem that lets squares compose into routes, together with the
projection-compatibility condition that pasting actually requires
(\S\ref{sec:pasting}); and the direction axis at work --- abstraction,
refinement, and what transfers (\S\ref{sec:lax}). Fidelity,
coverage, and the end-to-end guarantee build on this base in
\Cref{sec:fidelity}.

\subsection{Languages, behaviors, projections}
\label{sec:languages}

\begin{definition}[Language]
\label{def:language}
A \emph{language} $A$ is a triple $(\Prog{A}, F_A, \sem{\cdot}{A})$ where
$\Prog{A}$ is a decidable set of \emph{programs}; $F_A$ is a finite set of
named \emph{observable fields}, each $f \in F_A$ with a value domain
$V_f$; and $\sem{\cdot}{A} : \Prog{A} \rightharpoonup \Beh{A}$ is the
\emph{reference semantics}, a partial function into behaviors.
An \emph{observable state} is a total map
$\sigma \in \Obs{A}$, where $\Obs{A} = \prod_{f \in F_A} V_f$; a
\emph{behavior} is a finite sequence of observable states,
$\Beh{A} = \Obs{A}^{*}$, recorded \emph{post-step}: the $i$-th state is
the observable state after the $i$-th transition.
\end{definition}

The only admission criterion for a language is that $\sem{\cdot}{A}$ is
\emph{definable} --- an ISA manual, a logic's model theory, a reaction-network
semantics, and a molecular-graph notation all qualify. Nothing in this
section assumes languages are executable or even computational.

The definition also fixes the calculus's scope, and we state the
restrictions plainly: behaviors are \emph{finite} and the reference
semantics is a \emph{function}, so nondeterminism, divergence, and
reactive I/O are outside this calculus. Behavior \emph{inclusion}, by
contrast, is inside --- in exactly one controlled form: a square is a
\emph{directional} claim (\Cref{def:pair,def:faithful}), checked as a
projected trace \emph{equality} along a declared witness embedding. A
two-sided ($\mathsf{exact}$) square asserts equality of kept behaviors;
a one-sided ($\mathsf{over}$) square asserts the simulation half ---
every source behavior has a target counterpart, and the target may have
more; nothing else is admitted. What remains outside is unrestricted
refinement of nondeterministic semantics --- the reference semantics
itself stays a function. Projected trace equality along an embedding is
the right notion for the near-structure-preserving translations the
platform builds and checks, and the wrong one for optimizing
compilation, whose home is full simulation over retimed structure;
accordingly, the one genuinely optimizing hop in the instantiation
carries no square at all and is handled differently
(\S\ref{sec:composition}, \Cref{sec:instantiation}). Languages whose
official semantics is not a function (C, with its unspecified orders and
undefined behavior) enter through a determinized fragment, with the gap
policed differentially (\Cref{sec:instantiation}).

A program here is a \emph{closed} configuration: free inputs, when present,
are part of $p$ (a program together with an input valuation). Open programs
--- programs quantified over their inputs --- enter the square through
their \emph{closing valuations} (\Cref{def:faithful} quantifies over
them) and take center stage in \S\ref{sec:endtoend}, where a solver
decides over all inputs and a witness closes one.

\begin{definition}[Interpreter]
\label{def:interpreter}
An \emph{interpreter} for $A$ is a computable partial function
$I_A : \Prog{A} \rightharpoonup \Beh{A}$ that is \emph{pure}: identical
input yields byte-identical output, on every host (the contract;
\S\ref{sec:eval-perf} enforces it by twice-and-diff on the evaluation
host). Outside its domain
it fails with a typed $\unsupported$ verdict naming the construct;
$\dom(I_A)$ is $A$'s \emph{covered fragment}.
\end{definition}

\begin{assumption}[Interpreter adequacy]
\label{asm:adequacy}
$I_A(p) = \sem{p}{A}$ for all $p \in \dom(I_A)$.
\end{assumption}

Adequacy is an assumption, not a theorem, and we keep it visible: it is the
irreducible residue in every trusted-computing-base computation of
\S\ref{sec:tcb}. It is discharged \emph{empirically}, by differential
testing against an independently produced oracle for the same semantics
(the Sail-generated RISC-V simulator; \Cref{sec:instantiation} --- which
also notes the one honest exception, Python, whose pinned runtime is
both semantics and interpreter, leaving adequacy there with no
independent reference), and its failure modes are cross-checked by the
same machinery that checks translators.

\begin{definition}[Projection]
\label{def:projection}
A \emph{projection} over $A$ is a subset $\proj \subseteq F_A$. It lifts
to states by restriction, $\proj(\sigma) = \sigma\!\restriction_{\proj}$,
and to behaviors pointwise,
$\proj(\sigma_1 \cdots \sigma_n) = \proj(\sigma_1) \cdots \proj(\sigma_n)$
(in particular, projected behaviors of different lengths are unequal).
Write $b \eqproj{\proj} b'$ for $\proj(b) = \proj(b')$.
\end{definition}

\subsection{Pairs and the directional square}
\label{sec:pairs}

\begin{definition}[Pair]
\label{def:pair}
A \emph{pair} $P : A \to B$ is a tuple
$(A, B, T, \carry, \proj_P, d_P)$ where
$T : \Prog{A} \rightharpoonup \Prog{B}$ is the \emph{translator};
$\carry : \Beh{B} \rightharpoonup \Beh{A}$ is the
\emph{target-to-source interpreter}, which re-expresses a target behavior
as a source behavior; $\proj_P \subseteq F_A$ is the pair's
\emph{declared projection} --- exactly the source observables it promises
to preserve; and $d_P \in \{\mathsf{exact}, \mathsf{over}\}$ is its
\emph{declared direction}. An $\mathsf{over}$ (over-approximating) pair
additionally ships a pure \emph{witness embedding} $W$, mapping each
closing valuation $x$ of a source program $p$ to a closing valuation
$W(p,x)$ of $T(p)$; for an $\mathsf{exact}$ pair, $W$ is the identity.
$T$, $\carry$, and $W$ are pure (as in \Cref{def:interpreter});
$I_A$ and $I_B$ are owned by the languages and shared by every pair that
touches them. The pair's domain $\dom(P)$ is the set of $p$ with
$p \in \dom(I_A) \cap \dom(T)$, $T(p) \in \dom(I_B)$, and
$I_B(T(p)) \in \dom(\carry)$. Like the projection, the direction is a
declared, \emph{protected} part of the contract.
\end{definition}

\begin{definition}[Faithfulness: the directional square]
\label{def:faithful}
\label{def:lax}
$P$ is \emph{faithful at} $p \in \dom(P)$ when its square commutes at
$p$ \emph{along the witness embedding}: for every closing valuation $x$
of $p$,
\[
\proj_P\bigl(I_A(p(x))\bigr) \;=\;
\proj_P\bigl(\carry(I_B(\,T(p)(W(p,x))\,))\bigr),
\]
where a closed program has the single trivial valuation. For an
$\mathsf{exact}$ pair ($W = \mathrm{id}$) this is the familiar
commuting square $\proj_P(I_A(p)) = \proj_P(\carry(I_B(T(p))))$:
\[
\begin{tikzcd}[column sep=large]
p \arrow[r, "T"] \arrow[d, "I_A"'] & T(p) \arrow[d, "I_B"] \\
I_A(p) & \carry(I_B(T(p))) \arrow[l, "\eqproj{\proj_P}"']
\end{tikzcd}
\]
$P$ is faithful on $C \subseteq \dom(P)$ when it is faithful at every
$p \in C$.
\end{definition}

The direction is the \emph{reading} of that equality. The embedding
selects which target valuations the square checks: exactly those in its
range. When $W$ is the identity --- surjective onto the closing
valuations of $T(p)$ --- every target valuation is checked and the
target has no behaviors beyond the source's kept ones: the square is
two-sided, a bisimulation on the kept observables, and that is all
$d_P = \mathsf{exact}$ asserts. When $W$ is not surjective, the target
valuations outside its range are the pair's \emph{added} behaviors ---
their existence is the point, not a defect --- and the same checked
equality asserts only the simulation half,
$I_A(p) \sqsubseteq_{\proj_P} \carry(I_B(T(p)))$: every source behavior
has a target counterpart on the kept observables, and the target may
have more. Exactness is thus not a separate notion but the
identity-embedding special case of one square; \Cref{prop:exactid}
makes the equivalence formal, and it is machine-checked. Two
consequences are worth stating before any theory. First, the check
itself is indifferent to the direction --- run both legs, closed at $x$
and $W(p,x)$, compare under $\proj_P$ --- so everything below that
speaks of running, localizing, or ratcheting squares applies to both
directions verbatim; only what an agreement \emph{means} differs.
Second, that reading is why the direction is protected like the
projection: an $\mathsf{over}$ square passed off as $\mathsf{exact}$
would launder added behaviors into a meaning-preservation claim.

Three remarks. First, everything in \Cref{def:faithful} is computable:
faithfulness at a given $p$ is a \emph{decidable, executable check} --- run
both sides, compare under $\proj_P$. We call this check the
\emph{square oracle}; in fault-tolerance terms it is an acceptance test
in the sense of recovery blocks \citep{randell1975recovery}, turning a
silently wrong --- in effect Byzantine --- translator into a
\emph{fail-stop} component \citep{schlichting1983failstop}. Second, when the check fails it fails at a
\emph{first} position: a least step index $i$ and a field
$f \in \proj_P$ on which the two sides differ. The oracle reports
$(i, f)$ --- a translator (or interpreter, or carry-back) defect, localized
to a step and a named observable. Third, the projection is part of the
contract, not a tuning knob: $\proj_P$ states what ``preserves meaning''
promises for this pair, and $F_A \setminus \proj_P$ --- the pair's
\emph{loss} --- states what it does not.

The carry-back $\carry$ is what makes a target-side result \emph{mean}
something at the source: a solver counterexample or a target trace is
re-expressed as a source behavior. It is pair-specific because the
correspondence it encodes is pair-specific; it is also why the square's
oracle is as expressive as the translator it checks.

\begin{genexample}[One square, run --- and failing]
The \texttt{JAL} probe of the RV64IMC inventory
(\Cref{def:coverage}) is a two-instruction program whose taken jump
leaves the code image:
\begin{center}
\footnotesize\ttfamily
\begin{tabular}{@{}lll@{}}
0x0: & jal x1, 8 & ; link x1 := 4, jump to 0x8 \\
0x4: & ecall     & ; (jumped over) \\
\end{tabular}
\end{center}
The left leg runs the shared RISC-V interpreter $I_A$; the right leg
translates to BTOR2, interprets, and carries the trace back, compared
under $\proj \supseteq \{\texttt{pc}, \texttt{x1}, \texttt{halted}\}$.
The square is exact ($W = \mathrm{id}$) --- as is every registered
pair's but the abstraction endo-pairs' (\S\ref{sec:lax}).
Post-step states, with the right leg as the
\texttt{riscv-btor2} translator emitted them at version 0.1:
\begin{center}
\footnotesize
\begin{tabular}{@{}clll@{}}
\toprule
step & field & $I_A(p)$ & $\carry(I_B(T(p)))$ \\
\midrule
0 & \texttt{pc}, \texttt{x1}, \texttt{halted} & 8, 4, F & 8, 4, F \\
1 & \texttt{pc}, \texttt{x1} & 8, 4 & 8, 4 \\
  & \texttt{halted} & \textbf{T} & \textbf{F} \\
\bottomrule
\end{tabular}
\end{center}
Address \texttt{0x8} matches no instruction: the reference semantics
halts on the fetch miss, but the translator had modeled only the
instructions --- its transition system left \texttt{halted} low and the
machine stuck. The oracle reports the first divergence as
$(i, f) = (1, \texttt{halted})$ and, run hop by hop
(\Cref{cor:localization}), indicts the \texttt{riscv-btor2} hop. This
is incident I21 (\S\ref{sec:bugs}) exactly as the first conjoined
coverage run caught it --- on this probe, on every taken-jump probe,
and, the same morning, in two more independently written lowerings
with the same missing edge. The fix is a versioned translator bump
($0.1 \to 0.2$, an off-code$\to$halt transition) with the taken-jump
squares as permanent regressions; the all-pass rows of
\Cref{tab:capability} are made of exactly such checks.
\end{genexample}

\subsection{Pasting: when squares compose}
\label{sec:pasting}

Two pairs compose when the middle language is shared:
$P_1 = (A, B, T_1, \carry_1, \proj_1, d_1)$ and
$P_2 = (B, C, T_2, \carry_2, \proj_2, d_2)$ yield the candidate composite
\[
P_2 \circ P_1 \;=\; (A,\, C,\; T_2 \circ T_1,\; \carry_1 \circ \carry_2,\;
\proj,\; d_1 \wedge d_2)
\]
for a projection $\proj \subseteq \proj_1$ to be determined, with
direction the meet on the chain $\mathsf{exact} > \mathsf{over}$ --- a
composite is exact iff both hops are --- witness embedding the
composition $W_2 \circ W_1$ (an exact hop contributing the identity),
and $\dom(P_2 \circ P_1)$ the evident composite domain. Folklore says the
squares paste. They do not paste for free: $P_2$ guarantees agreement of
the two $B$-behaviors only \emph{up to} $\proj_2$, and $\carry_1$ consumes
whole $B$-behaviors. If $\carry_1$ reads a $B$-field that $P_2$ discards,
the outer rectangle can fail while both inner squares hold. The missing
condition is that $\carry_1$ must not look outside what $P_2$ preserves:

\begin{definition}[Support]
\label{def:support}
$\carry_1$ is \emph{$(\proj_2 \Rightarrow \proj)$-supported} when for all
$b, b' \in \dom(\carry_1)$:
$\proj_2(b) = \proj_2(b')$ implies
$\proj(\carry_1(b)) = \proj(\carry_1(b'))$.
Equivalently: $\proj \circ \carry_1$ descends to a well-defined map on
$\proj_2$-equivalence classes --- a congruence condition, matching the
mechanization exactly.
\end{definition}

\begin{theorem}[Pasting]
\label{thm:pasting}
Let $p \in \dom(P_2 \circ P_1)$, let $P_1$ be faithful at $p$, let $P_2$
be faithful at $T_1(p)$, let $\proj \subseteq \proj_1$, and let
$\carry_1$ be $(\proj_2 \Rightarrow \proj)$-supported. Then the
composite $P_2 \circ P_1$ is faithful at $p$ with respect to $\proj$,
along the composed embedding $W_2 \circ W_1$ and with direction
$d_1 \wedge d_2$; for the exact--exact instance, writing
$r = T_2(T_1(p))$,
\[
\proj(I_A(p)) \;=\; \proj\bigl(\carry_1(\carry_2(I_C(r)))\bigr).
\]
\end{theorem}

\begin{proof}
Write $q = T_1(p)$. Then $\proj(I_A(p)) = \proj(\carry_1(I_B(q)))$ since
$\proj \subseteq \proj_1$ and $P_1$ is faithful at $p$. Faithfulness of
$P_2$ at $q$ gives
$\proj_2(I_B(q)) = \proj_2(\carry_2(I_C(r)))$, so the
support condition applies to the two $B$-behaviors and yields
$\proj(\carry_1(I_B(q))) = \proj(\carry_1(\carry_2(I_C(r))))$.
Chain the equalities. For the directional form, run the same chain per
closing valuation, the legs closed at $x$, $W_1(p,x)$, and
$W_2(T_1(p), W_1(p,x))$; mechanized as \texttt{lax\_pasting} and,
telescoped over routes with the direction meet,
\texttt{DRoute.lax\_route\_pasting} and
\texttt{DRoute.direction\_exact\_iff} (\S\ref{sec:mech}).
(Details: \Cref{app:proofs}.)
\end{proof}

The support condition is not decorative. Take $F_B = \{g, h\}$ with
$\proj_2 = \{g\}$, and let $\carry_1$ copy field $h$ into a
$\proj_1$-kept source field: $P_2$ can be faithful --- it preserves
$g$ --- while $\carry_2(I_C(r))$ differs from $I_B(q)$ on $h$, which
$\proj_2$ permits. The outer rectangle then fails on the copied field
although both inner squares commute.

\paragraph{Composed keep and loss}
In the implementable special case --- $\carry_1$ \emph{fieldwise} and
step-aligned, each source field $f$ computed from a dependency set
$\mathit{deps}(f) \subseteq F_B$ of target fields --- the support
condition has a syntactic reading, and the largest $\proj$ satisfying
\Cref{thm:pasting} is
\[
\keep(P_2 \circ P_1) \;=\;
  \{\, f \in \proj_1 \mid \mathit{deps}(f) \subseteq \proj_2 \,\},
\]
a source field survives the route iff $P_1$ keeps it and it is computed
only from target fields $P_2$ keeps. Dually
$\lost(P_2 \circ P_1) = F_A \setminus \keep(P_2 \circ P_1)$: route loss is
the union of per-hop losses, each pulled back along the carry-backs. This
is the formal content of the platform rule that a route must \emph{declare
its cumulative loss}: the observables a destination can still speak about
are exactly those no hop discarded. (When $\carry_1$ changes step
granularity --- one source step spanning several target steps, as in
C-to-ISA --- the same statement holds with $\mathit{deps}$ taken over the
spanned window; appendix~A.2.)

\begin{corollary}[Localization]
\label{cor:localization}
Under the support condition, if the composite square fails at $p$ (both
sides defined), then $P_1$'s square fails at $p$ or $P_2$'s square fails
at $T_1(p)$. Inductively, along a route
$R = P_n \circ \cdots \circ P_1$ satisfying the pairwise support
conditions, a failing outer rectangle implies a failing inner square at
some hop $i$ --- and running the square oracle hop by hop finds the
least such $i$, then the least failing step and field. (The check is
direction-indifferent, so this holds along any directional route ---
divergences localize the same way whether the failing square's reading
was $\equiv$ or $\sqsubseteq$.)
\end{corollary}

\Cref{cor:localization} is the debugging story of the whole platform: a
divergence anywhere along a route is never a mystery about the route; it
is an indictment of one hop, one step, one observable. It is also what
disagreement between \emph{routes} will lean on in \S\ref{sec:branching}.

\paragraph{Routes}
The registry of languages and pairs forms a directed multigraph; a
\emph{route} $R : A \to Z$ is a path through it, with composed translator
$T_R$, composed carry-back $\carry_R$, composed projection $\keep(R)$,
composed direction and embedding as above, and
faithfulness given by iterating \Cref{thm:pasting}. Purity composes too,
and gives caching:

\begin{proposition}[Determinism and caching]
\label{prop:cache}
If every component function of every hop of $R$ is pure, then $T_R$ and
$\carry_R$ are pure, and a content-addressed cache keyed by the input
hash and the component versions returns, on a hit, exactly what
recomputation would return --- across the whole route. If some hop is
impure, re-running the route on the same input and
diffing bytes at every hop (\emph{twice-and-diff}) can detect it and
localize it to the first hop whose bytes differ --- a sound but
incomplete detector: two coinciding runs of an impure hop pass.
\end{proposition}

\subsection{The direction axis: abstraction as a pair}
\label{sec:lax}

Every square so far shown ran exact, and most registered pairs do. The
$\mathsf{over}$ half of \Cref{def:faithful} earns its place on cost:
an \emph{abstraction} maps a program to a smaller or cheaper model with
\emph{more} behaviors --- cut a state variable's update and let it
range freely, and every source behavior survives while new, spurious
ones appear. Such a translation is exactly what makes an expensive
universal question tractable, and in this calculus it is not a special
mechanism bolted onto a solver run but an ordinary pair whose declared
direction is $\mathsf{over}$: same oracle, same coverage conjunction
(\Cref{def:coverage}), same ratchet (\Cref{prop:ratchet}), same
negative controls, checked along its witness embedding. Two results
pin the axis down: exactness is a special case, and universal verdicts
transfer.

\begin{proposition}[Exactness is the identity embedding]
\label{prop:exactid}
Under the specialization obligation of \Cref{thm:universal}(iv) ---
translating the closed instance agrees with closing the one open
translation --- a pair with $W = \mathrm{id}$ is faithful at an open
$p$ in the sense of \Cref{def:faithful} iff its two-sided square
commutes at every closed instance $p(x)$. The exact calculus of
\S\ref{sec:pairs}--\S\ref{sec:pasting} is thus the
$W = \mathrm{id}$ fragment of the directional one --- literally: the
mechanization derives each side from the other
(\texttt{laxFaithful\_id\_iff\_faithful}, axiom-free;
\S\ref{sec:mech}).
\end{proposition}

\begin{proof}
With $W = \mathrm{id}$ the two legs close at the same valuation, and
the specialization obligation identifies the closed translation
$T(p(x))$ with the closed instance $T(p)(x)$ of the open artifact;
the quantification over valuations is then the quantification over
closed instances. \qed
\end{proof}

\begin{proposition}[Universal transfer across over-approximation]
\label{prop:lax}
Let $R : A \to Z$ have composed direction $\mathsf{over}$ and be
faithful (\Cref{def:faithful}) on a fragment containing $p$, with
$\varphi$ over $\keep(R)$. If no valuation of $T_R(p)$ exhibits
$\varphi$ within the declared bounds, then no valuation of $p$ does:
\emph{universal verdicts transfer across over-approximation}. An
existential verdict does not transfer; it remains subject to the
carried-back replay of \Cref{thm:existential}, whose failure on an
$\mathsf{over}$ route has a second, benign reading --- a \emph{spurious
counterexample}, a demand to refine the abstraction rather than an
indictment of a hop.
\end{proposition}

\begin{proof}
Contrapositive: a valuation $x$ with $\varphi$ observed in
$\sem{p(x)}{A}$ yields, by faithfulness along the composed embedding,
the valuation $W_R(p,x)$ of $T_R(p)$ exhibiting $\varphi$ on the kept
observables --- present within bounds, contradicting the premise.
(Mechanized: \texttt{lax\_universal\_transfer}, \S\ref{sec:mech}.)
\end{proof}

Together with \Cref{thm:existential,thm:universal} this closes a loop
the platform otherwise leaves to external tools:
counterexample-guided abstraction refinement. Ask the universal
question on the abstraction; a transferred \textsc{unreachable} is an
answer about the source (\Cref{prop:lax}), a \textsc{reachable}
either replays at the source (\Cref{thm:existential} --- a true
counterexample) or fails to --- a spurious one, and the refinement the
failure demands is not solver-internal state but a \emph{new
registered pair} with a smaller havoc set, graded and reusable like
any other. The artifact instantiates this with an over-approximating
\emph{endo-pair} (source and target language coincide):
\texttt{btor2-havoc}, localization abstraction on the bit-level hub,
whose lax square, negative controls, and refinement loop run in the
repository's gate (a second $\mathsf{over}$ endo-pair, interval
abstraction, is registered as a brief).

\paragraph{The dual, for symmetry}
The axis has an evident third value the platform deliberately does not
yet inhabit: an \emph{under}-approximating square --- the embedding
running the other way, the target with \emph{fewer} behaviors ---
where existential verdicts would transfer and universal ones would
not, the mirror image of \Cref{prop:lax}. Nothing in this paper
depends on it; we note it because the directional reading makes the
symmetry visible, and its verdict-transfer discipline would fall out
of the same monotonicity.

The mechanization of \S\ref{sec:mech} covers the directional calculus
in full: the directional square, pasting along the composed embedding
(binary and telescoped) with the direction meet, and both propositions
above are machine-checked (\texttt{Calculus/Lax.lean}) with the same
axiom footprint as the exact fragment.
